\documentclass[main.tex]{subfiles}

\begin{document}

\section{Rewards Pass-Through and Network Incidence}\label{sec:appendix-network-incidence}

In this appendix we make a narrow point within the original two-sided framework of \citet{Rochet.Tirole2003}: when the network prices optimally, the rate at which interchange passes through to rewards is orthogonal to whether the network bears the incidence of an interchange-fee cut.
A natural intuition holds that if networks do not fully pass interchange through to rewards, they must retain part of any cut and so absorb some of its incidence.
The envelope theorem refutes this: because the network was already balancing merchant acceptance against rewards optimally, a small forced reduction in the merchant fee changes its profit only at second order, no matter how aggressively it re-optimizes rewards.
The pass-through rate governs how the cut is split between cardholders and merchants; it says nothing about the network's own (zero) first-order incidence.

\subsection{Setup}\label{subsec:netinc-setup}

The monopoly network is a combined entity (network, issuers, acquirers) that sets two per-transaction prices: $p_s$ charged to merchants and $p_b$ charged to cardholders.
Rewards correspond to $p_b<0$.

Payment volumes depend on both consumer and merchant adoption. We model this with a standard two-sided demand structure.
A share $D_b(p_b)$ of consumers adopt, and a share $D_s(p_s)$ of merchants accept the card.
Card transactions only occur when both consumers want to use the card and merchants accept, and so total payment volume equals $Q=D_b D_s$.
Profit is
\begin{equation}
\pi(p_b,p_s)\;=\;(p_b+p_s-c)\,D_b(p_b)\,D_s(p_s).
\label{eq:netinc-profit}
\end{equation}

In the unconstrained equilibrium, the network sets optimal prices $(p_b^*,p_s^*)$ to maximize~\eqref{eq:netinc-profit}.
Writing the markup as $\phi\equiv p_b^*+p_s^*-c$, the interior FOCs $\partial\pi/\partial p_b=D_b D_s+\phi\,D_b'D_s=0$ and $\partial\pi/\partial p_s=D_b D_s+\phi\,D_b D_s'=0$ yield the slope identities
\begin{equation}
D_b\;=\;-\phi\,D_b',\qquad D_s\;=\;-\phi\,D_s'\qquad\text{at }(p_b^*,p_s^*).
\label{eq:netinc-foc-identities}
\end{equation}

\subsection{The Network's Reward Response}\label{subsec:netinc-envelope}

Suppose a regulator caps $p_s\leq p_s^*-\varepsilon$ for $\varepsilon>0$ small.
The network re-optimizes $p_b$ subject to the constrained merchant price, yielding $p_b^*(\varepsilon)$ defined by the residual buyer-side FOC
\begin{equation}
D_b(p_b^*)\;+\;(p_b^*+p_s^*-\varepsilon-c)\,D_b'(p_b^*)\;=\;0.
\label{eq:netinc-constrained-foc}
\end{equation}
Implicitly differentiating~\eqref{eq:netinc-constrained-foc} at $\varepsilon=0$ pins down the cardholder-side response,
\begin{equation}
\rho^R\;\equiv\;\frac{dp_b^*}{d\varepsilon}\bigg|_{\varepsilon=0}\;=\;\frac{D_b'(p_b^*)}{2D_b'(p_b^*)+\phi\,D_b''(p_b^*)}.
\label{eq:netinc-alpha}
\end{equation}
$\rho^R$ measures how much of a merchant-fee cut the network claws back by raising the cardholder price (i.e., cutting rewards).
We use the same symbol $\rho^R$ as in Appendix~\ref{sec:appendix-theory}: both measure the fraction of a merchant-fee change that reaches cardholders.
There the parameter enters a reduced-form rebate rule; here it is microfounded as the network's profit-maximizing response.

Under log-concave demand, $\rho^R$ lies strictly between zero and one, so the network generically offsets only part of a merchant-fee cut with lower rewards.
Log-concavity captures most common demand curves, including linear and isoelastic, and is standard in the literature on imperfect pass-through \citep{Weyl.Fabinger2013}.

\begin{lemma}[Range of $\rho^R$]\label{lem:alpha-range}
Under log-concavity of cardholder demand at the monopoly point, $\rho^R\in(0,1)$.
\end{lemma}

\begin{proof}
The denominator in~\eqref{eq:netinc-alpha} equals $\partial^2\pi/\partial p_b^2\,/\,D_s$, which is negative by the unconstrained second-order condition. With $D_b'<0$, this gives $\rho^R>0$. The upper bound $\rho^R<1$ requires $D_b'+\phi\,D_b''<0$; substituting the FOC identity $D_b=-\phi\,D_b'$ rewrites this as $D_b\,D_b''\leq(D_b')^2$, equivalently $(\log D_b)''\leq 0$.
\end{proof}

\subsection{Zero Network Incidence}

That the network re-optimizes rewards is precisely why the cap leaves its profit untouched to first order. The network was already balancing merchant acceptance against cardholder rewards optimally, so a small forced perturbation has no first-order profit effect by the envelope theorem.

\begin{theorem}[Zero Network Incidence]\label{thm:network-incidence}
Suppose the network's pre-cap prices satisfy the unconstrained FOCs with negative-definite Hessian and $D_b$ is log-concave at $p_b^*$. Then $d\pi^*/d\varepsilon\big|_{\varepsilon=0}=0$.
\end{theorem}

\begin{proof}
Let $\pi^*(\varepsilon)\equiv\pi(p_b^*(\varepsilon),p_s^*-\varepsilon)$. Differentiating,
\[
\frac{d\pi^*}{d\varepsilon}\;=\;\frac{\partial\pi}{\partial p_b}\bigg|_*\cdot\frac{dp_b^*}{d\varepsilon}\;+\;\frac{\partial\pi}{\partial p_s}\bigg|_*\cdot(-1).
\]
At $\varepsilon=0$ both partials vanish by the unconstrained FOCs, so $d\pi^*/d\varepsilon|_{\varepsilon=0}=0$. \qed
\end{proof}

The pass-through rate $\rho^R$ appears nowhere in this result, and the reason is visible in the derivative itself.
Writing $d\pi^*/d\varepsilon=(\partial\pi/\partial p_b)\,\rho^R-\partial\pi/\partial p_s$, the reward response enters only as the coefficient on the buyer-side margin $\partial\pi/\partial p_b$---a margin the network has already driven to zero.
However aggressively the network claws back rewards, that response multiplies a vanishing margin and cannot move first-order profit.
Equivalently, $\rho^R$ is a curvature object: \eqref{eq:netinc-alpha} depends on $D_b''$, a feature of the second-order condition, whereas network incidence is a statement about the gradient, which is zero by construction at the optimum.
The extent of rewards pass-through and the incidence of the cap are different questions about the same optimum.
The local result also survives in estimated structural models: even large caps deliver only modest changes in network profits \citep{Wang2025}.

To leading order the cut therefore redistributes between the network's two sides without touching the network itself: cardholders lose $\rho^R\varepsilon$ per inframarginal transaction (the clawed-back rewards), merchants gain the full $\varepsilon$, and $\Delta_{\text{network}}=0$.
The pass-through rate fixes how this surplus divides between cardholders and merchants---the redistribution studied in Appendix~\ref{sec:appendix-theory}, which carries the same accounting through to consumers via retail pass-through $\rho^P$---but it never reaches the network.
The orthogonality is therefore specific to network incidence: $\rho^R$ remains first-order for cardholder welfare and for the efficiency gain the cap delivers, and the network's slice is the one quantity it leaves untouched.

\end{document}
